% !TEX root = ../main.tex
\chapter{Introduction; logic of propositions and predicates}\label{ch:logic}
\begin{paracol}{2}

\begin{sources}
\begin{tabularx}{\linewidth}{@{}l l L@{}}
\toprule
\textbf{Video} & \textbf{Length} & \textbf{Content}\\
\midrule
\seg{1}{1}{Lecture 1: Introduction / Logic of propositions and predicates} & 1:40:48 & Aim and structure of the
course; propositional logic (unary and binary operators, implication, proof by contradiction, NAND);
predicate logic (quantifiers, order of quantification); axiomatic systems, proofs, consistency;
consistency of propositional logic; G\"odel's incompleteness theorem (statement and idea); outlook on
axiomatic set theory.\\
\bottomrule
\end{tabularx}

\smallskip
\textbf{Transcript} \file{materials/transcripts/C01-V01.txt} (YouTube automatic captions).
Recorded at FAU Erlangen, winter semester 2013/14.
\end{sources}

\begin{reading}[Where this lecture sits in the other notes]
None of the other Schuller courses starts from logic: \emph{Gravity and Light} (Lecture~1) and
\emph{Theoretische Mechanik} (Lecture~2) start directly from topology, which here is Lecture~4.
The overview of the ``building'' in \cref{sec:lo-structure} is the same picture as the bird's-eye view of
\emph{Theoretische Mechanik}, Lecture~1. Links to the chapters will be added when those notes exist.
\end{reading}

% ==================================================================
\section{What the course is for}\label{sec:lo-aims}

Theoretical physics casts our concepts about the real world into rigorous mathematical form, ``for better or
worse''; it does so not for its own sake, but to explore fully what those concepts imply \ts{1}{1}{0:07}.
Schuller puts the spirit of the subject in the words of Wittgenstein: \emph{what we cannot speak about clearly,
we must pass over in silence} \ts{1}{1}{0:51}. If a concept about the world cannot be put into precise
mathematical form, this usually shows that some aspect of the concept has not been understood.
\begin{coursenote}[The two Wittgenstein quotes]
The first is the last proposition (7) of the \emph{Tractatus Logico-Philosophicus} (1921): ``Wovon man nicht
sprechen kann, dar\"uber muss man schweigen.'' The second is close to Tractatus 5.43, where it is said of the
propositions of \emph{logic}: ``Alle S\"atze der Logik sagen aber dasselbe. N\"amlich nichts.''
\end{coursenote}

Mathematics, however, is \emph{just a language}. To extract physical conclusions from a mathematical
formulation one must interpret it, and that is not the task of mathematics but of physicists, which is where
it becomes difficult again \ts{1}{1}{1:32}. Again with Wittgenstein: the theorems of mathematics all say the
same thing, namely nothing \ts{1}{1}{1:52}. This does not mean that mathematics is useless. A theorem of the
form ``$A$ if and only if $B$'' says nothing that $A$ did not already say: logically it is a tautology. But
psychologically, for our understanding of $A$, a reformulation in terms of $B$ may be very useful
\ts{1}{1}{2:13}.

With this understanding, the idea of the course is to provide the proper language for theoretical physics
\ts{1}{1}{2:56}.

\begin{intuition}[The aim in one line]
Provide proper mathematical language for classical mechanics, electromagnetism, quantum mechanics and
statistical physics \ts{1}{1}{3:19}.
\end{intuition}

These are the subjects of the modules TP1 to TP4 at Erlangen \ts{1}{1}{4:04}. The course will not revise all
the mathematics they need; it develops the mathematics from a higher point of view, now that some knowledge
of these subjects can be assumed. So far they have been studied with a fair amount of \emph{analysis}, of
\emph{algebra} (mostly linear, sometimes non-linear) and of \emph{geometry} in one form or another
\ts{1}{1}{4:25}: classical mechanics and electromagnetism appeal to geometric intuition to set up and solve a
problem; quantum mechanics, in its modern form, is done with linear algebra; statistical physics uses some
geometry, maybe without noticing, when it looks at convex cones of states \ts{1}{1}{5:06}.

Drawing a circle around each of the three fields, the physical theories sit roughly in the intersections
\ts{1}{1}{5:26}, as in \cref{fig:lo-venn}. Classical mechanics and electrodynamics use analysis and geometry,
and so does general relativity; quantum mechanics lies in the intersection of analysis and algebra, usually
called \emph{functional analysis}; statistical physics, ``to do it justice'', lies more or less between algebra
and geometry \ts{1}{1}{5:51}. This is not a pure picture, only where the intersections come into play. In the
middle, where all three intersect, sits \emph{differential geometry}, and it is mainly differential geometry
that the course is about \ts{1}{1}{6:33}.

\begin{nfigure}
\begin{tikzpicture}[scale=0.95]
  \draw[thick,defcol] (0,0) circle (2.1) node[above left=1.05cm and 0.6cm] {\textbf{Analysis}};
  \draw[thick,thmcol] (2.6,0) circle (2.1) node[above right=1.05cm and 0.6cm] {\textbf{Algebra}};
  \draw[thick,intcol] (1.3,-2.2) circle (2.1) node[below=1.35cm] {\textbf{Geometry}};
  \node[font=\footnotesize,align=center] at (1.3,1.25) {QM\\(functional\\analysis)};
  \node[font=\footnotesize,align=center] at (-0.45,-1.45) {CM, EM,\\GR};
  \node[font=\footnotesize,align=center] at (3.05,-1.45) {statistical\\physics};
  \node[font=\footnotesize\bfseries,align=center] at (1.3,-0.65) {differential\\geometry};
\end{tikzpicture}
\caption{Where the physical theories sit among the three mathematical fields \ts{1}{1}{5:51}. The course is
about the centre.}\label{fig:lo-venn}
\end{nfigure}

% ==================================================================
\section{The structure of the course: a building}\label{sec:lo-structure}

The development of differential geometry starts with sets \ts{1}{1}{7:17}. Every space of physics (the
configuration space or the phase space of classical mechanics, the space in which electromagnetism works, the
physical space over which the Hilbert space of quantum mechanics is built) is, at its coarsest level, a set of
points \ts{1}{1}{7:38}. But what precisely is a set? Saying ``a collection of elements'' only raises the
questions of what a collection is and what elements are. There is a fundamental problem: if we start writing a
book on mathematics whose pages are still empty, what can the first definition be? A definition needs notions
we already have \ts{1}{1}{7:58}. The way out is to start \emph{axiomatically}; but then, in what language are
the axioms written? So below set theory there is one more building block: \emph{logic}, propositional and
predicate logic, which defines our language \ts{1}{1}{8:40}. The building then goes up as follows
(\cref{fig:lo-building}).

\begin{enumerate}
  \item \textbf{Logic and axiomatic set theory}: the language and the axioms for sets \ts{1}{1}{9:02}.
  \item \textbf{Topology.} The four subjects need a notion of \emph{continuity}: famously, particles do not
        jump; a trajectory does not stop somewhere and start again somewhere else. Continuity is obtained by
        distinguishing certain subsets of the basic set, which leads to topological spaces \ts{1}{1}{9:23}.
  \item \textbf{Topological manifolds.} There are so many topological spaces that they cannot even be
        classified, so we restrict to those that are immediately interesting for physics (topologists may
        disagree that they are the only interesting ones): spaces that locally around each point look like
        some $\R^d$ \ts{1}{1}{10:04}.
  \item \textbf{Differentiable manifolds.} Classical physics needs derivatives, for instance the velocity of a
        curve, which cannot be defined on a topological manifold. Differentiable manifolds are ``an entirely
        new beast'' and need the technology of \emph{charts}, which makes contact with the formalism of all
        the standard courses: these have all been formulated in charts so far \ts{1}{1}{10:46}.
  \item \textbf{Bundles}: to each point of a manifold one attaches another manifold, for instance a vector
        space at each point (the tangent space) \ts{1}{1}{12:52}.
  \item \textbf{Geometry}: tensor fields on these structures \ts{1}{1}{14:58}.
  \item \textbf{Physics} on top: classical mechanics, electromagnetism, quantum mechanics, statistical physics,
        special and general relativity \ts{1}{1}{16:02}.
\end{enumerate}

\begin{intuition}[Charts and atlases]
A chart is what a sailor uses: travelling across the globe, you take out the atlas, turn to the chart of the
place where you are and navigate in it. The chart must be a one-to-one picture of the world locally. If you
move so that you would leave the page, you turn to another page that slightly overlaps with the first and
navigate on \ts{1}{1}{11:30}. Looking for overlapping charts and understanding how the overlaps work is the
whole philosophy of manifolds, and of differentiable manifolds in particular: the charts are pieces of
$\R^d$, where differentiation is understood, and the notion of differentiation is lifted from the charts to
the manifold \ts{1}{1}{12:11}.
\end{intuition}

Bundles are a very important concept in differential geometry, and much of modern physics can be cast in
their language: particle physics, gravity, and certainly quantum mechanics \ts{1}{1}{13:14}. Two warnings come
with them \ts{1}{1}{13:35}:
\begin{itemize}
  \item there is \emph{no such thing as a position vector}, however often it is repeated: a position vector
        is just a collection of coordinates, and these do not constitute a vector;
  \item there is \emph{no such thing as a wave function} in quantum mechanics: it is not a function at all,
        but a section of a complex line bundle over physical space \ts{1}{1}{14:16}.
\end{itemize}
One gets away with both notions as long as one does not look too closely; then they fail. ``You cannot
understand this now, but you will at the end of the course'' \ts{1}{1}{14:37}.

The geometry comes as tensor fields on these structures \ts{1}{1}{14:58}: a \emph{symplectic} geometry, the
geometry of the phase space of classical mechanics, and in fact of virtually all dynamical systems of interest
\ts{1}{1}{15:19}; a \emph{metric} geometry, needed to some extent in electromagnetism (there, really, only a
volume form is needed), and certainly in special and general relativity, which need a so-called
\emph{Lorentzian} metric \ts{1}{1}{15:41}. Up to the geometry everything is ``physical mathematics''; above it
is physics, and all of it can be properly understood by mastering this mathematics \ts{1}{1}{16:23}. The
course develops the mathematics and then shows, mainly by examples, how the insights play out at the level of
physics \ts{1}{1}{16:43}. The picture should not be taken too seriously (bundles, for instance, sometimes come
before or alongside differentiable manifolds) \ts{1}{1}{14:58}.

\begin{nfigure}
\begin{tikzpicture}[every node/.style={draw,minimum width=9.5cm,minimum height=0.62cm,font=\small,align=center}]
  \node[fill=thmcol!10] (ph) at (0,6.2) {\textbf{Physics}: CM, EM, QM, statistical physics, SR, GR};
  \node[fill=intcol!10] (ge) at (0,5.3) {\textbf{Geometry}: tensor fields (symplectic, metric, Lorentzian, volume form)};
  \node[fill=defcol!8]  (bu) at (0,4.4) {\textbf{Bundles} (no position vectors, no wave functions)};
  \node[fill=defcol!8]  (dm) at (0,3.5) {\textbf{Differentiable manifolds} (charts, atlases)};
  \node[fill=defcol!8]  (tm) at (0,2.6) {\textbf{Topological manifolds} (locally $\R^d$)};
  \node[fill=defcol!8]  (to) at (0,1.7) {\textbf{Topological spaces} (continuity)};
  \node[fill=srccol!10] (se) at (0,0.8) {\textbf{Axiomatic set theory}};
  \node[fill=srccol!20] (lo) at (0,-0.1) {\textbf{Logic}: propositional and predicate logic};
\end{tikzpicture}
\caption{The building of the course, from the bottom up \ts{1}{1}{8:40}. Up to geometry it is mathematics;
the top floor is physics.}\label{fig:lo-building}
\end{nfigure}

% ==================================================================
\section{Propositional logic}\label{sec:lo-prop}

Logic is treated as the first part of Chapter~1, \emph{Axiomatic set theory}, because it is the language
in which set theory will be written \ts{1}{1}{17:39}.

\begin{definition}[New concepts: proposition]\label{def:lo-proposition}
A \textbf{proposition} $p$ is a variable that can take the values \emph{true} or \emph{false}, and no others
\ts{1}{1}{18:32}.
\end{definition}
The name comes from the Latin \emph{propositio}, ``a setting forth'', a statement put forward.

That is all a proposition is for propositional logic. In particular, it is not its task to decide whether a
complex statement such as ``there is extraterrestrial life'' is true or not: propositional logic starts from
complete propositions and only assumes that each is true or false \ts{1}{1}{19:14}. Nor is it its task to
decide whether a sentence such as ``in winter it is colder than outside'' is a proposition at all (this one does
not seem particularly reasonable) \ts{1}{1}{19:56}. Propositions are its atomic, elementary notion
\ts{1}{1}{20:16}.

New propositions are built from given ones by \textbf{logical operators} \ts{1}{1}{20:38}.

\paragraph{Unary operators.} A unary operator takes one proposition $p$ and makes a new proposition from it
\ts{1}{1}{21:05}. Since $p$ can only be true or false, there are exactly four of them \ts{1}{1}{21:47}: the
negation $\neg p$, the identity $\mathrm{id}\,p$, the tautology $\top p$ (always true, whatever $p$ is) and the
contradiction $\bot p$ (always false) \ts{1}{1}{22:07}:
\begin{center}
\renewcommand{\arraystretch}{1.35}
\begin{tabular}{c|cccc}
$p$ & $\neg p$ & $\mathrm{id}\,p$ & $\top p$ & $\bot p$\\\hline
T & F & T & T & F\\
F & T & F & T & F
\end{tabular}
\end{center}
These are all the possibilities \ts{1}{1}{22:28}.

\paragraph{Binary operators.} A binary operator takes two given propositions $p,q$; the pair can take the
four values TT, TF, FT, FF \ts{1}{1}{22:48}. For each of the four rows we may choose T or F, so there are
$2^4=16$ binary operators \ts{1}{1}{23:29}. Some of them:
\begin{center}
\renewcommand{\arraystretch}{1.35}
\begin{tabular}{cc|ccccc}
$p$ & $q$ & $p\wedge q$ & $p\vee q$ & $p\veebar q$ & $p\Rightarrow q$ & $p\Leftrightarrow q$\\\hline
T & T & T & T & F & T & T\\
T & F & F & T & T & F & F\\
F & T & F & T & T & T & F\\
F & F & F & F & F & T & T
\end{tabular}
\end{center}
with $\wedge$ ``and'' \ts{1}{1}{23:09}, $\vee$ ``or'' (inclusive) and $\veebar$ the exclusive or
\ts{1}{1}{25:33}.

\begin{coursenote}[Why not ``maps'' from $\{T,F\}\times\{T,F\}$?]
Asked whether a binary operator is a map from the Cartesian product of $\{\text{true},\text{false}\}$ with
itself to $\{\text{true},\text{false}\}$, Schuller answers that this is the idea, but he cannot say it that way:
we do not yet know what a set or a Cartesian product is \ts{1}{1}{24:10}. Even more abstractly, one can
write down \emph{deduction rules} and never write a table. He does not, because for standard logic it would
serve no purpose \ts{1}{1}{24:31}. For non-standard logics, such as intuitionist logic, one removes some
deduction rules and gets a logic that can no longer be expressed by tables: ``very esoteric'', very
interesting, but every course has to define its bottom line, and this one presents tables instead of
deduction rules and axioms, which is sufficient for what follows \ts{1}{1}{24:52}.
\end{coursenote}

\begin{history}[``Or'': inclusive or exclusive?]
As told in class \ts{1}{1}{25:54}: in the Middle Ages people discussed whether the ``or'' in some sentence meant
the inclusive or the exclusive one, and appealed to the Bible to see how God used it. Around the turn of the
20th century Frege, Wittgenstein and others said: it is just a table, a definition; there are 16 binary
operators, each of which may or may not deserve a name.\par
\smallskip
\emph{References.} The medieval anecdote is reported as said in class (not checked). Truth tables in their
modern form: L.~Wittgenstein, \emph{Tractatus Logico-Philosophicus} (1921), 4.31 and 5.101; E.~L.~Post,
``Introduction to a general theory of elementary propositions'', \emph{Amer.\ J.\ Math.} 43 (1921) 163--185.
G.~Frege, \emph{Begriffsschrift} (1879).
\end{history}

This is all ``child's play'', but one binary operator is sometimes ill understood: the \textbf{implication}
$\Rightarrow$ \ts{1}{1}{26:35}. From school we use the arrow to mean ``I wrote down something true, did something
valid, and arrived here''. But the implication is a binary operator like $\wedge$: it takes two propositions
and makes one new proposition, $p\Rightarrow q$, which as a whole is true or false \ts{1}{1}{26:57}. Its table
is the column above: from true follows true, true; from true follows false, false (``you don't pass the
test''); from false follows true, true; and from false follows false, also true \ts{1}{1}{27:18}. This is the
definition \ts{1}{1}{28:01}.

\begin{definition}[New concepts: ex falso quodlibet]\label{def:lo-exfalso}
The principle that from a false assumption anything can be concluded, the implication as a whole being true,
is called \textbf{ex falso quodlibet} (Latin: ``from the false, whatever you like'') \ts{1}{1}{28:22}.
\end{definition}

The first two rows are intuitive; why define the last two like this? The answer is hidden in a little
theorem \ts{1}{1}{28:42}. To state it we need the \textbf{equivalence} $p\Leftrightarrow q$, true when both
propositions have the same value and false otherwise (table above) \ts{1}{1}{29:05}.

\begin{theorem}[Contraposition]\label{thm:lo-contra}
$(p\Rightarrow q)\ \Leftrightarrow\ \bigl((\neg q)\Rightarrow(\neg p)\bigr)$ \ts{1}{1}{29:25}.
\end{theorem}
\begin{proof}
Write out the truth table in full \ts{1}{1}{32:10}:
\begin{center}
\renewcommand{\arraystretch}{1.3}
\begin{tabular}{cc|cc|cc}
$p$ & $q$ & $\neg p$ & $\neg q$ & $p\Rightarrow q$ & $\neg q\Rightarrow\neg p$\\\hline
T & T & F & F & T & T\\
T & F & F & T & F & F\\
F & T & T & F & T & T\\
F & F & T & T & T & T
\end{tabular}
\end{center}
The last two columns agree in every row \ts{1}{1}{33:12}, so the equivalence is true whatever $p$ and $q$ are.
\end{proof}
Note that putting $\neg$ in front exchanges the order of $p$ and $q$; since binding strengths have not been
fixed yet, brackets are used \ts{1}{1}{29:46}. Proving such a statement by writing out the truth table is ``the
first vanilla problem in any exam'' \ts{1}{1}{32:10}.

\begin{corollary}[Proof by contradiction]\label{cor:lo-contradiction}
Assertions can be proven by way of contradiction \ts{1}{1}{30:06}: to prove that $q$ follows from the true
assumption $p$, one may equivalently assume that $q$ is false and show that then $p$ cannot be true.
\end{corollary}
This is a very elementary principle and a very sharp tool of mathematics \ts{1}{1}{30:48}.

\begin{secondary}[Intuitionist logic]
If logic is formulated by deduction rules and axioms instead of tables, one can ``surgically'' remove a rule
so that \cref{thm:lo-contra} no longer holds; the result cannot be written with tables any more, and it gives a
mathematics in which proofs by contradiction are not allowed \ts{1}{1}{30:48}. In every good mathematics
department there is one professor who follows this \emph{intuitionist} logic, and ``it's a pain'': everything
must be proven constructively \ts{1}{1}{31:29}. Some make the point that a constructive result is more
meaningful, and an even deeper one: that proofs by contradiction cannot be trusted. The course does not go into
this: its logic allows proofs by contradiction, essentially by the definition of $\Rightarrow$
\ts{1}{1}{31:49}.
\end{secondary}

\begin{exercise}[A tautology]\label{ex:lo-taut}
Write the truth table of $p\vee\neg p$ and of $(p\wedge(p\Rightarrow q))\Rightarrow q$. Which of the two is
true for every choice of $p,q$?
\end{exercise}

\begin{remark}[Binding strength]\label{rem:lo-binding}
We agree on the decreasing binding strength $\neg,\ \wedge,\ \vee,\ \Rightarrow,\ \Leftrightarrow$: $\neg$ binds
most strongly, $\Leftrightarrow$ least \ts{1}{1}{33:37}. With more symbols the sequence must be extended
\ts{1}{1}{33:59}.
\end{remark}

\begin{remark}[Higher-order operators and NAND]\label{rem:lo-nand}
An operator $\heartsuit(p_1,\dots,p_n)$ that takes $n$ propositions is an \textbf{operator of order $n$}: unary
operators have order 1, binary ones order 2 \ts{1}{1}{34:41}. All operators of any order can be constructed from
one single binary operator \ts{1}{1}{35:02}, the \textbf{NAND} operator (``not and''), $p\uparrow q:\Leftrightarrow
\neg(p\wedge q)$ \ts{1}{1}{35:24}:
\begin{center}
\renewcommand{\arraystretch}{1.35}
\begin{tabular}{cc|c}
$p$ & $q$ & $p\uparrow q$\\\hline
T&T&F\\ T&F&T\\ F&T&T\\ F&F&T
\end{tabular}
\end{center}
So the binding strengths are not really needed: they are a convenience \ts{1}{1}{36:07}.
\end{remark}

\begin{exercise}[First steps with NAND]\label{ex:lo-nand}
Express $\neg p$ and $p\wedge q$ using only $\uparrow$ (``finger practice'', as on the first problem sheet
\ts{1}{1}{36:27}). The full statement is \cref{ex:lo-nand-full} in the appendix.
\end{exercise}

Asked again about the equivalence operator: it is true exactly when both propositions take the same value;
that is its definition \ts{1}{1}{36:48}.

% ==================================================================
\section{Predicate logic}\label{sec:lo-pred}

\begin{definition}[New concepts: predicate]\label{def:lo-predicate}
A \textbf{predicate} is, informally, a proposition-valued function of some variable or variables
\ts{1}{1}{37:38}. For example $P(x)$ is true or false depending on $x$ \ts{1}{1}{37:58}; $Q(x,y)$ is true or false
depending on the combination of $x$ and $y$ \ts{1}{1}{38:21}.
\end{definition}
The name comes from the Latin \emph{praedicare}, ``to proclaim'': what is asserted of a subject.

As propositional logic does not study how complex propositions arise, predicate logic, strictly speaking, does
not say how predicates are built \ts{1}{1}{38:41}. One could think of $Q(x,y)$ being true if $x>y$ for real
numbers $x,y$; but constructing such a predicate needs a further language for combining objects $x$ and $y$
\ts{1}{1}{39:02}. And from which set do $x$ and $y$ come? \emph{No such thing}: what $x$ and $y$ are, where they
come from and what their nature is, is left entirely open; they are just variables \ts{1}{1}{39:44}. This is
weird, because since elementary school we have been conditioned always to say which set a variable comes from;
but it must be left open, because the notion of set is to be defined later \emph{in this language}
\ts{1}{1}{40:05}.

\begin{coursenote}[Outlook: the element predicate]
Set theory will introduce a fundamental predicate of two variables, $x\in y$, which for two objects is either
true or false \ts{1}{1}{40:25}. If $x$ and $y$ are sets, $x\in y$ is a predicate of two variables. This is only
an outlook: in predicate logic the symbol $\in$ does not exist yet \ts{1}{1}{41:08}. ``It is always important to
know what a subject is not talking about'': here, how predicates are built \ts{1}{1}{41:28}.
\end{coursenote}

\paragraph{New predicates from old.} One can build new predicates from given ones \ts{1}{1}{41:48}, for instance
\[
  Q(x,y,z):\Leftrightarrow P(x)\wedge R(y,z),
\]
where $:\Leftrightarrow$ means ``is defined to be equivalent to'' (the analogue of $:=$): once $x,y,z$ are
supplied, $P(x)$ and $R(y,z)$ are propositions, which can be joined by $\wedge$ into a new predicate, now of three
variables \ts{1}{1}{42:10}.

\paragraph{Quantifiers.} More interesting: a predicate of one variable can be converted into a proposition
\ts{1}{1}{42:51}.

\begin{definition}[New concepts: universal quantifier]\label{def:lo-forall}
For a predicate $P$ of one variable, $\forall x\colon P(x)$ is the proposition, read ``for all $x$, $P(x)$ is
true'' \ts{1}{1}{44:14}, which is \emph{defined} to be true if $P(x)$ is true independently of $x$, that is,
whatever is substituted for $x$ \ts{1}{1}{44:38}.
\end{definition}
The whole expression is a proposition constructed from a predicate of one variable \ts{1}{1}{43:54}.

\begin{example}[A ``feel-good'' predicate]
Predicate logic does not provide specific predicates, but as a ``feel-good example'' (and a feel-good example
always involves some imprecision) \ts{1}{1}{45:26}: let
$P(x):\Leftrightarrow (x\text{ is a human being})\Rightarrow(x\text{ has been created})$, which ``includes all
faiths'' \ts{1}{1}{46:10}. The whole proposition is the implication; where the $x$ come from is hidden in its
first part. This is certainly always true, so $\forall x\colon P(x)$ is true \ts{1}{1}{46:36}.
\end{example}

A student asks whether a predicate can be made into a proposition in other ways, ``for some $x$''
\ts{1}{1}{46:57}. Indeed: the brother of the universal quantifier \ts{1}{1}{47:18}.

\begin{definition}[New concepts: existential quantifier]\label{def:lo-exists}
$\exists x\colon P(x)$, read ``there exists an $x$ such that $P(x)$'' \ts{1}{1}{47:59}, is the proposition
\[
  \exists x\colon P(x)\ :\Leftrightarrow\ \neg\bigl(\forall x\colon \neg P(x)\bigr),
\]
the negated universal quantification of the negated predicate \ts{1}{1}{48:19}.
\end{definition}

\begin{corollary}\label{cor:lo-neg}
$\forall x\colon\neg P(x)\ \Leftrightarrow\ \neg\bigl(\exists x\colon P(x)\bigr)$ \ts{1}{1}{48:40}: if something is
not true for any $x$, then it is not true that there exists an $x$ for which it is true \ts{1}{1}{49:03}.
\end{corollary}

\begin{aside}[The hairdresser's door]
Simple once written down, but careful with negations and quantifiers \ts{1}{1}{49:28}. On the doors of many
barber shops it says: ``What hairdressers can do, only hairdressers can do.'' It sounds good, but it means that
hairdressers can do nothing a non-hairdresser can do: if a non-hairdresser can draw nice pictures, a hairdresser
cannot \ts{1}{1}{49:48}.
\end{aside}

\begin{remark}[Quantifying predicates of several variables]\label{rem:lo-free}
Quantification also applies to predicates of more than one variable \ts{1}{1}{50:30}. $\forall x\colon P(x,y)$ is a
predicate of the one variable $y$ that was not quantified \ts{1}{1}{51:14}: $x$ is the \textbf{bound} variable,
$y$, which survives the quantification, the \textbf{free} variable \ts{1}{1}{51:35}. The remaining variable can
be quantified again \ts{1}{1}{51:56}.
\end{remark}

\begin{remark}[The order of quantification matters]\label{rem:lo-order}
In general \ts{1}{1}{52:17}
\[
  \forall x\ \exists y\colon P(x,y)\qquad\text{and}\qquad \exists y\ \forall x\colon P(x,y)
\]
are different propositions \ts{1}{1}{52:37}. One should therefore seriously avoid, although it is tricky in the
heat of a calculation, writing the quantifiers \emph{after} the expression (``\dots for all $x,y,z$''): when
written informally behind, it is unclear where they should be pulled to. Write all quantifiers in front, in
order \ts{1}{1}{52:58}.
\end{remark}

\begin{example}[Neutral and inverse elements]
For the additive neutral element of the real numbers: \emph{there exists} $0\in\R$ such that \emph{for all}
$x\in\R$, $x+0=x$; the same $0$ works for every $x$ \ts{1}{1}{53:40}. For inverses: \emph{for every} $x$ \emph{there
exists} a $y$ with $x+y=0$; $y$ may be chosen depending on $x$ \ts{1}{1}{54:01}. (Schuller first stated the
neutral element the wrong way round, and corrected himself: ``you see, already a mistake again''.)
\end{example}

\begin{exercise}[Negating quantifiers]\label{ex:lo-negq}
Using \cref{def:lo-exists}, show that $\neg(\forall x\ \exists y\colon P(x,y))\Leftrightarrow\exists x\ \forall y\colon
\neg P(x,y)$. Then negate the statement ``every real number has an additive inverse''.
\end{exercise}

% ==================================================================
\section{Axiomatic systems and the theory of proofs}\label{sec:lo-proofs}

Proofs are met in physics and, above all, in mathematics lectures; but what \emph{is} a proof
\ts{1}{1}{55:17}? A way of arguing why something is true given that something else was assumed true; but what are
the valid rules for writing one down? This can be defined, and there are subtleties which make some proofs
that are given really non-proofs \ts{1}{1}{55:38}. All of this prepares the axiomatic system of set theory
\ts{1}{1}{55:58}.

\begin{definition}[New concepts: axiomatic system]\label{def:lo-axsys}
An \textbf{axiomatic system} is a finite sequence of propositions $a_1,a_2,\dots,a_N$, which are called the
\textbf{axioms} \ts{1}{1}{56:19}\ts{1}{1}{58:27}.
\end{definition}

\begin{aside}[Pre-mathematical numbers]
A smart objection: what do ``finite'' and the numbers $1,2,\dots,N$ mean before sets are defined? One could
write $a_{/},a_{//},a_{///},\dots$ with strokes instead; logicians call these \emph{pre-mathematical numbers}:
``it's what you're doing in Bavarian pubs'' (counting beers by strokes on the beer mat) \ts{1}{1}{56:44}.
\end{aside}

\begin{coursenote}[Disclaimer]
A mathematical logician in the audience would shout ``this is a little too na\"ive'': the course does not use the
full formal approach of first-order predicate logic (this already showed in the use of tables) \ts{1}{1}{57:24}.
It gives the right idea, conceptually exactly what logicians do, without the formal nitty-gritty, which for
what follows would be fully redundant. It is not what a professional logician would write in a research
monograph, but nothing of it is wrong \ts{1}{1}{57:46}.
\end{coursenote}

The idea: first the axioms are given, then a homework problem ``prove $p$ starting from the axioms''
\ts{1}{1}{59:08}.

\begin{definition}[New concepts: proof]\label{def:lo-proof}
A \textbf{proof} of a proposition $p$ within an axiomatic system $a_1,\dots,a_N$ is a finite sequence of
propositions (or propositional schemes) $q_1,\dots,q_M$ (the steps of the proof) \ts{1}{1}{59:30} such that
$q_M=p$ (the last step is the proposition to be proven) \ts{1}{1}{1:00:15} and, for every $1\le j\le M$, one of the
following holds \ts{1}{1}{1:02:41}\ts{1}{1}{1:03:03}:
\begin{itemize}
  \item[(A)] $q_j$ is a proposition from the list of axioms \ts{1}{1}{1:03:23};
  \item[(T)] $q_j$ is a tautology \ts{1}{1}{1:03:44};
  \item[(M)] there exist $1\le m,n<j$ such that $(q_m\wedge q_n)\Rightarrow q_j$ is true \ts{1}{1}{1:05:28}\ts{1}{1}{1:06:10}.
\end{itemize}
The letters stand for \emph{axiom}, \emph{tautology} and \emph{modus ponens} \ts{1}{1}{1:05:28}.
\end{definition}
So at any step one may write down an axiom (A), a plain tautology (T), or a proposition deduced from two
\emph{previous} steps (M) \ts{1}{1}{1:07:12}. \emph{Modus ponens} (Latin, ``the mode that affirms'') is a very old
deduction rule.

\begin{definition}[New concepts: tautology]\label{def:lo-tautology}
A \textbf{tautology} is a proposition that is true independently of the elementary propositions it is made of
\ts{1}{1}{1:04:05}; for example $p\vee\neg p$. Generically the value of an expression containing $p$ depends on $p$,
but this particular combination is always true, as the table shows \ts{1}{1}{1:04:25}.
\end{definition}

\begin{remark}[Propositional schemes]
One single axiom can have the form ``$\forall x\colon P(x)$''. It is just one proposition, but one may then use
the individual instances $P(x)$, which are true in maybe infinitely many cases; this is called a
\textbf{propositional scheme}. Still, the steps of a proof must be \emph{finitely many}, and this is very
important \ts{1}{1}{1:00:38}\ts{1}{1}{1:01:18}.
\end{remark}

\begin{aside}[``This proof is not finite'']
In his first semester Schuller was, as he says today, fortunate that his analysis professor was an intuitionist:
no proofs by contradiction, and he was sometimes really mad at textbook proofs of analysis theorems, often saying
``this proof is not finite'' [?: the transcript has ``not cyanide'', probably German \emph{nicht finit}]. The
students did not understand what he meant; with the formal theory of proofs, where a proof is a finite sequence
of propositions, they did \ts{1}{1}{1:01:39}.
\end{aside}

\paragraph{Is rule (M) needed?} A student asks whether (M) does not already follow from (A) and (T)
\ts{1}{1}{1:07:54}. It is needed, ``definitely, seriously, desperately, fully'': the intuition ``if this is true,
then I can do this'' is kindergarten intuition; $\Rightarrow$ is just a binary operator, with no built-in
direction \ts{1}{1}{1:08:15}. There is not yet enough technology for a meaningful example, but the course will
write out \emph{one} proof in full according to this scheme: the uniqueness of the empty set. An axiom will
stipulate that there exists an empty set; showing that there is only one is intuitively clear, but written out in
full it takes about 16 lines \ts{1}{1}{1:08:36}\ts{1}{1}{1:08:56}.

\begin{intuition}[Machine code and C++]
All other proofs are given the usual way, by quick arguments. Not even the most celebrated theorems in the best
journals are proven in the formal fashion, which would take ``millions and billions'' of lines for comparatively
simple statements \ts{1}{1}{1:09:37}. The formal proof is like machine code, zeros and ones; the proofs we write
are in a higher language, like C++ \ts{1}{1}{1:09:58}. The understanding is that a written proof is a proof only if
fellow mathematicians agree that it could in principle be broken down into a proof of the formal type
\ts{1}{1}{1:10:18}. If somebody suspects that it would not (for instance, not into finitely many steps), they object;
and as long as somebody says ``I'm not convinced'', the proof has to be broken down further and further
\ts{1}{1}{1:10:39}. ``Don't tell anyone: you can obstruct any lecture in mathematics'' this way \ts{1}{1}{1:10:59}.
\end{intuition}

Why may (M) use any two previous steps, even with others in between, but no later ones? Everything before the
current step is an axiom or a tautology or was reached from them by implication arrows; going beyond the current
step would amount to using equivalence arrows (``and later I show\dots''), and apart from circularity it is the
directedness of the implication that allows exactly the previous steps \ts{1}{1}{1:11:21}\ts{1}{1}{1:11:41}.

\begin{remark}[Notation]
If $p$ can be proven from the axiomatic system $a_1,\dots,a_N$, one writes $a_1,\dots,a_N\vdash p$ (``proves
$p$'') \ts{1}{1}{1:12:23}.
\end{remark}

\begin{remark}[Recognising versus finding]
The definition makes it easy to \emph{recognise} a proof: a computer can check whether a given sequence is a proof
of $p$ \ts{1}{1}{1:13:05}\ts{1}{1}{1:13:28}. \emph{Finding} a proof is an altogether different matter: even for set
problems, which can be solved, it is not easy, and sometimes it is so hard that whoever finds a proof first
becomes a celebrated mathematician \ts{1}{1}{1:13:49}. The theory of proofs gives no hint on how to prove things;
it tells you whether what you have written is a proof, which can be very useful \ts{1}{1}{1:14:30}.
\end{remark}

\begin{remark}[Tautologies are redundant axioms]\label{rem:lo-tautax}
A tautology among the axioms can be removed from the list without impairing the power of the axiomatic system:
the same propositions can still be proven, because the tautology can be pulled into any proof by (T)
\ts{1}{1}{1:14:56}\ts{1}{1}{1:15:41}. The extreme case: the axiomatic system of \emph{propositional logic} is the
\emph{empty sequence}, since everything one wants to show there is a tautology \ts{1}{1}{1:16:22}\ts{1}{1}{1:16:46}.
\end{remark}

\begin{definition}[New concepts: consistency]\label{def:lo-consistent}
An axiomatic system $a_1,\dots,a_N$ is \textbf{consistent} if there exists a proposition $q$ that cannot be proven
from it \ts{1}{1}{1:17:34}\ts{1}{1}{1:17:57}:
\[
  \exists q\colon\ \neg\,(a_1,\dots,a_N\vdash q).
\]
\end{definition}

The idea behind this maybe surprising definition \ts{1}{1}{1:18:40}: consider an axiomatic system containing
contradictory propositions, say $s$ and, later in the list, $\neg s$ \ts{1}{1}{1:19:28}\ts{1}{1}{1:19:54}. By (A) both
$s$ and $\neg s$ can be pulled into a proof; then by (M), for an arbitrary $q$, it suffices that
$(s\wedge\neg s)\Rightarrow q$ is true \ts{1}{1}{1:20:20}. It is a tautology: $s\wedge\neg s$ is always false, and by
the definition of the implication arrow, ex falso quodlibet, the whole statement is always true
\ts{1}{1}{1:21:02}\ts{1}{1}{1:21:23}. So any statement $q$ can be proven, and also $\neg q$ \ts{1}{1}{1:21:44}. Not being
able to prove every statement is a sign of not having contradictory axioms; hence the definition excludes exactly
this \ts{1}{1}{1:22:05}.

\begin{theorem}\label{thm:lo-propconsistent}
Propositional logic is consistent \ts{1}{1}{1:22:45}.
\end{theorem}
\begin{proof}
It suffices to exhibit a proposition that cannot be proven within propositional logic \ts{1}{1}{1:23:11}. Its
axiomatic system is the empty sequence (\cref{rem:lo-tautax}), so only (T) and (M) can be used in a proof
\ts{1}{1}{1:23:56}. The only way to introduce a new step is to write down a tautology, and (M) only allows, from
previous tautologies, a $q_j$ for which the whole implication is a tautology; so only tautologies can be proven
\ts{1}{1}{1:24:17}\ts{1}{1}{1:24:43}. But $q\wedge\neg q$ is not a tautology, so it cannot be proven
\ts{1}{1}{1:25:03}. Hence propositional logic is consistent \ts{1}{1}{1:25:26}.
\end{proof}
This is really trivial \ts{1}{1}{1:25:46}. The important remark is that while it is perfectly clear how to
\emph{define} consistency, it is perfectly difficult to \emph{prove} it for a given axiomatic system, propositional
logic being a big exception. For a more powerful system, like the axioms of set theory, it is very difficult; and
famous results, first obtained by G\"odel, say that under certain circumstances it is not even possible
\ts{1}{1}{1:26:07}\ts{1}{1}{1:26:27}. Only a rough taste of one of them \ts{1}{1}{1:26:48}:

\begin{theorem}[G\"odel, roughly and non-technically]\label{thm:lo-goedel}
Any axiomatic system that is powerful enough to encode the elementary arithmetic of the natural numbers is either
inconsistent or contains a proposition that can be neither proven nor disproven \ts{1}{1}{1:27:09}\ts{1}{1}{1:28:19}.
\end{theorem}
``Elementary arithmetic'' is imprecise language: roughly elementary-school mathematics (prime numbers, products);
what exactly is needed one sees from the proof \ts{1}{1}{1:27:57}. Variants of the formulation say that the
system contains true statements that cannot be proven \ts{1}{1}{1:29:10}.

\begin{history}[G\"odel's incompleteness theorems]
This sent shockwaves through the mathematical world \ts{1}{1}{1:29:10}: in mathematics, of all the arts, the notion
of truth had seemed clear: something was true if it could be proven from a system of assumptions
\ts{1}{1}{1:29:32}. Then G\"odel came along and, ironically, \emph{proved} this and variants of it
\ts{1}{1}{1:29:52}.\par
\smallskip
\emph{References.} K.~G\"odel, ``\"Uber formal unentscheidbare S\"atze der Principia Mathematica und verwandter
Systeme I'', \emph{Monatshefte f\"ur Mathematik und Physik} 38 (1931) 173--198 (the first and second
incompleteness theorems; announced at K\"onigsberg in September 1930).
\end{history}

\begin{secondary}[The idea of G\"odel's proof]
The proof is involved; its basic idea is the following \ts{1}{1}{1:29:52}.
\begin{enumerate}
  \item Assign to each mathematical, and also \emph{meta}-mathematical, statement a number, now called its
        \textbf{G\"odel number} \ts{1}{1}{1:30:15}. In a mathematical statement such as $a^2+b^2=c^2$, or one with
        quantifiers, roughly the plus sign gets a number, the exponent gets one, every variable a different one,
        the equal sign another \ts{1}{1}{1:31:04}. One can then talk \emph{about} the equation: ``the left-hand
        side has more symbols than the right-hand side'' is a meta-mathematical statement, and G\"odel showed how
        to assign a number even to such statements \ts{1}{1}{1:31:24}\ts{1}{1}{1:31:44}.
  \item The construction: every elementary symbol is assigned a prime; one takes powers of primes and multiplies
        them, getting big numbers, huge ones for meta-mathematical statements \ts{1}{1}{1:32:07}. Since numbers are
        multiplied, elementary arithmetic is needed: this explains the mysterious condition on the axiomatic
        system \ts{1}{1}{1:32:29}.
  \item Then a barber-type argument (``the barber shaves all people in his village who do not shave themselves'')
        identifies a proposition that is neither provable nor disprovable \ts{1}{1}{1:32:50}\ts{1}{1}{1:33:14}.
\end{enumerate}
The construction is reflexive, like Cantor's proof that there are uncountably many real numbers
\ts{1}{1}{1:34:08}: one assumes that every number written as an integer, a decimal point and decimal places
$d_1d_2d_3\dots$ can be listed, and uses the numbering to construct a new real number that is not in the list
\ts{1}{1}{1:34:29}. In G\"odel's proof one assigns G\"odel numbers to statements and produces a statement that uses
the G\"odel number of the very same statement, which leads to the contradiction \ts{1}{1}{1:34:49}.
\end{secondary}

With the definitions of this lecture (axiomatic system, provable, hence disprovable: the negation can be proven,
and consistent) the \emph{statement} of the theorem is fully understood; the proof is a different matter
\ts{1}{1}{1:35:10}.

% ==================================================================
\section{Outlook: axiomatic set theory}\label{sec:lo-outlook}

This completes the language in which axiomatic set theory will be defined, starting next time
\ts{1}{1}{1:35:30}.
\begin{itemize}
  \item The \textbf{element symbol} $\in$ is nothing but a predicate of two variables, written in the middle
        ($x\in y$ instead of ${\in}(x,y)$) to save two brackets and a comma, as is often done for predicates of
        two variables \ts{1}{1}{1:35:51}\ts{1}{1}{1:36:11}.
  \item It is the \emph{only} elementary predicate of set theory: even the equality of sets will be defined in
        terms of it \ts{1}{1}{1:36:32}. Nobody will explain what it means that $x$ is an element of $y$; and both
        $x$ and $y$ will be sets: there is no hierarchy of ``grown-up sets'' and mere elements, sets can be
        contained in sets \ts{1}{1}{1:36:53}.
  \item A priori there is no reason to exclude a set being an element of itself \ts{1}{1}{1:37:13}. If one
        restricts to the sets that do not contain themselves, one proves in two lines that the collection of
        all such sets is not a set: asking whether it contains itself, both answers lead to a contradiction
        \ts{1}{1}{1:37:34}\ts{1}{1}{1:37:54}. So na\"ive set theory is tricky, and axioms are needed to regulate the
        use of the element predicate, a kind of \emph{operational} definition \ts{1}{1}{1:38:15}.
  \item Some axioms say that certain sets exist: an empty set; for a given set, its power set, the set of all its
        subsets \ts{1}{1}{1:38:35}. In fact the only set one really needs to postulate is the empty set: every set
        is ultimately constructed from it (the set containing the empty set, and so on) \ts{1}{1}{1:38:56}.
  \item There are about nine axioms (the number varies with how they are written) that reproduce the set theory
        on which all of modern mathematics is built \ts{1}{1}{1:39:16}.
  \item One of them is the \textbf{axiom of choice} (Zorn's lemma is an equivalent form). Why be so picky? Without
        it one cannot prove that every vector space has a basis, which is used very often
        \ts{1}{1}{1:39:16}\ts{1}{1}{1:39:37}.
\end{itemize}
Already at the bottom of the building there are screws that must be justified for the whole building to stand as
we want; the purpose of this lecture and the next is to look at the foundations and to have a justifiably good
feeling about the constructions to come \ts{1}{1}{1:39:58}\ts{1}{1}{1:40:18}.

\begin{coursenote}[Names used]
The paradox in the third point is Russell's paradox (B.~Russell, 1901; letter to Frege, 1902); the axioms are those
of Zermelo--Fraenkel set theory with choice (ZFC); both names come in the next lecture.
\end{coursenote}

% ==================================================================
\section*{Key concepts}
\begin{keyconcepts}
\textbf{Retrieval questions.}
\begin{enumerate}
  \item Why is $(F\Rightarrow F)$ true, and which theorem justifies this definition?
  \item How many unary and binary operators are there, and from which single operator can all be built?
  \item Define $\exists$ in terms of $\forall$. Why does the order of quantifiers matter?
  \item State the three rules (A), (T), (M) for a step of a proof.
  \item What does it mean for an axiomatic system to be consistent, and why does a contradiction among the axioms
        destroy it?
  \item Why is propositional logic consistent? What does G\"odel's theorem say?
\end{enumerate}
\textbf{Concepts.}
\begin{itemize}
  \item Proposition: a variable taking the values true or false (\cref{def:lo-proposition}).
  \item Implication: $p\Rightarrow q$ false only for $p$ true, $q$ false; ex falso quodlibet (\cref{def:lo-exfalso}).
  \item Contraposition: $(p\Rightarrow q)\Leftrightarrow(\neg q\Rightarrow\neg p)$; proof by contradiction
        (\cref{thm:lo-contra}, \cref{cor:lo-contradiction}).
  \item NAND: $p\uparrow q:\Leftrightarrow\neg(p\wedge q)$ generates all operators (\cref{rem:lo-nand}).
  \item Predicate; $\forall x\colon P(x)$; $\exists x\colon P(x):\Leftrightarrow\neg\forall x\colon\neg P(x)$
        (\cref{def:lo-forall,def:lo-exists}); bound and free variables.
  \item Axiomatic system, proof by (A), (T), (M), $a_1,\dots,a_N\vdash p$ (\cref{def:lo-axsys,def:lo-proof}).
  \item Consistency: $\exists q\colon\neg(a_1,\dots,a_N\vdash q)$ (\cref{def:lo-consistent}); propositional logic is
        consistent (\cref{thm:lo-propconsistent}); G\"odel (\cref{thm:lo-goedel}).
\end{itemize}
\end{keyconcepts}

\flushside
\end{paracol}
